\title{Controlling Text-to-Image Diffusion by Orthogonal Finetuning}

\begin{abstract}
Text-to-image diffusion models have demonstrated remarkable proficiency in synthesizing photorealistic images conditioned on textual descriptions. However, steering or constraining these powerful models to accomplish diverse downstream objectives remains a significant research challenge. In response, we present a systematic approach—Orthogonal Finetuning (OFT)—to adapt text-to-image diffusion models for specialized tasks. Distinct from prior techniques, OFT is able to rigorously maintain hyperspherical energy, a property reflecting the relational structure among neurons on the unit hypersphere. This preservation proves essential for maintaining the semantic fidelity of generations produced by text-to-image diffusion models. To further promote reliable and robust finetuning, we introduce Constrained Orthogonal Finetuning (COFT), which incorporates an explicit constraint on the hypersphere’s radius. Our study focuses on two representative finetuning applications: subject-driven generation, where the model learns to synthesize personalized images based on a limited set of subject images and a text cue; and controllable generation, where the model is equipped to interpret supplementary control signals alongside a text prompt. Through extensive experiments, we demonstrate that our OFT method not only surpasses existing strategies in terms of output quality, but also accelerates training convergence.
\end{abstract}

\section{Introduction}

State-of-the-art text-to-image diffusion models~\cite{saharia2022photorealistic,ramesh2022hierarchical,rombach2022high} have achieved notable success in generating high-quality images steered by text prompts. Despite these advances, relying exclusively on textual instructions can often result in ambiguous or imprecise image manipulation, lacking detailed and exacting control. To tackle these limitations, we focus on two specific text-to-image synthesis scenarios in this work:

\begin{itemize}[leftmargin=*,nosep]

    \item \textbf{Subject-driven generation}~\cite{ruiz2023dreambooth}: This setting involves producing images featuring the same subject in varied contexts, leveraging only a handful of images of the subject and a guiding text description.
    \item \textbf{Controllable generation}~\cite{zhang2023adding,mou2023t2i}: Here, the objective is to condition image creation not just on text, but also on auxiliary signals such as canny edges or segmentation masks, enabling the model to align generated content with both the text and additional controls.
\end{itemize}

The central challenge in both tasks is to adapt text-to-image diffusion models through finetuning while maintaining the generative quality achieved during pretraining. We identify two main criteria for an effective finetuning approach: (1) \emph{training efficiency}, defined by minimal increases in trainable parameters and training epochs, and (2) \emph{generalizability preservation}, referring to the retention of high-quality and diverse outputs. Current approaches to finetuning typically involve either incremental updates to existing neuron weights with a low learning rate (\emph{e.g.}, \cite{ruiz2023dreambooth}), or augmenting the architecture with minor, re-parameterized modules (\emph{e.g.}, \cite{hulora2022,zhang2023adding}). While straightforward, these strategies fall short of ensuring that the generative capacity acquired during pretraining is maintained. Moreover, a fundamental framework for developing robust finetuning strategies and for the principled selection of hyperparameters such as training epoch count remains absent. One notable obstacle lies in the absence of an objective metric to evaluate the retention of pretraining generative abilities. Finetuning methods commonly presuppose that smaller Euclidean distances between the parameters of finetuned and pretrained models correlate with better preservation of generative functions. Consequently, small learning rates are often employed. However, this reliance on Euclidean proximity is insufficient to encapsulate semantic retention. We contend that more sophisticated, structurally informed metrics are necessary to adequately assess and enhance the semantic consistency between the finetuned and original models, thereby supporting both the conservation of pretraining performance and the robustness of the finetuning process.

Drawing on findings that hyperspherical similarity effectively represents semantic relationships~\cite{liu2017deep,liu2018decoupled,chen2020angular}, we leverage hyperspherical energy~\cite{liu2018learning} as a means to describe the pairwise associations among neurons within a layer. Here, hyperspherical energy is calculated as the aggregate of hyperspherical similarities (such as cosine similarity) for all neuron pairs, thereby measuring the degree of their uniform distribution across the unit hypersphere~\cite{liu2021learning}. Our working assumption is that optimal finetuning should minimally disrupt the hyperspherical energy compared to the original pretrained model. A straightforward solution is to employ a regularization term to maintain constant hyperspherical energy throughout finetuning; however, this approach does not necessarily guarantee a small change in hyperspherical energy. As an alternative, we utilize a key invariance characteristic of hyperspherical energy—specifically, that the pairwise hyperspherical similarity remains intact under a global orthogonal transformation of the neuron vectors. This insight leads us to develop Orthogonal Finetuning~(OFT), an approach for tailoring large text-to-image diffusion models to downstream tasks while preserving hyperspherical energy. OFT centers on learning a shared orthogonal transformation at each layer, which maintains pairwise angular relationships among neurons. Alternatively, OFT may be interpreted as redefining the canonical coordinate axes for neurons within a layer. In addition, acknowledging that a reduced Euclidean distance between pretrained and finetuned models correlates with better retention of pretraining capabilities, we introduce Constrained Orthogonal Finetuning~(COFT)—an extension of OFT that restricts the finetuned model to lie on a hypersphere of fixed radius centered at the pretrained neuron positions.

The rationale behind utilizing orthogonal transformation for neuron finetuning is partially informed by the 2D Fourier transform, where an image can be separated into its magnitude and phase spectra. The phase spectrum encodes angular relationships between the input and the basis and is primarily responsible for preserving semantic information. For instance, if one were to combine the phase spectrum of an original image with an arbitrary magnitude spectrum, the essential semantic content of the image can still be retained upon reconstruction. This observation implies that modifying neuron orientations plays a pivotal role in altering the semantics of generated images—a core aim in both subject-driven and controllable generation. Nevertheless, altering neuron orientations without restraint risks compromising the generative abilities learned during pretraining. To prevent excessive freedom in this modification, we suggest maintaining the angles between all neuron pairs. This approach is grounded in the assumption that these inter-neuron angles encapsulate important knowledge within neural networks. Guided by this reasoning, we propose learning a layer-wise, shared orthogonal transformation that ensures the hyperspherical energy of each layer remains invariant.

Our approach is also motivated by prior work on orthogonal over-parameterized training~\cite{liu2021orthogonal}, where classification neural networks are trained from scratch through orthogonal transformations applied to random initializations. Randomly initialized neural networks are expected to possess low hyperspherical energy, and the primary objective of \cite{liu2021orthogonal} is to sustain this low energy throughout training, as reduced hyperspherical energy has been linked to improved classification generalization~\cite{liu2018learning,lin2020regularizing}. The findings in~\cite{liu2021orthogonal} demonstrate that orthogonal transformations are sufficiently expressive for constructing neural networks with strong generalization capabilities for classification. By contrast, our focus is on finetuning text-to-image diffusion models to achieve greater controllability and enhanced generative outcomes on downstream tasks. The divergence between OFT and the methodology of \cite{liu2021orthogonal} lies primarily in two areas. First, unlike~\cite{liu2021orthogonal}, which is geared towards minimizing hyperspherical energy, OFT aims to preserve the pretrained hyperspherical energy to retain the intrinsic structure of the original model during finetuning. In the context of diffusion model finetuning, reducing hyperspherical energy might compromise essential semantic structures. Second, OFT introduces an additional constraint by penalizing deviation from the pretrained model, resulting in a constrained version. This restriction is absent in~\cite{liu2021orthogonal}. The challenge of finetuning, therefore, is to effectively balance adaptability with retention of the original model's stability, a trade-off we contend is well addressed by our OFT framework. Our key contributions are summarized below:

\begin{itemize}[leftmargin=*,nosep]

    \item We introduce Orthogonal Finetuning (OFT), an innovative finetuning approach designed to enhance the controllability of text-to-image diffusion models. To address stability concerns, we also develop a constrained version that restricts angular deviation from the original pretrained parameters.
    \item Distinct from traditional finetuning methods, OFT enables model adaptation while theoretically ensuring that the hyperspherical energy is maintained. Our empirical analysis suggests that this quantity is critical for preserving the generative semantic properties of the pretrained network.
    \item We deploy OFT on two distinct challenges: subject-driven generation and controllable generation. Through extensive empirical evaluation, we show that OFT markedly outperforms prior methods with respect to image quality, training efficiency, and robustness during finetuning. Additionally, OFT is more sample-efficient, achieving satisfactory convergence even when limited training images and epochs are available.
    \item For evaluating controllable generation, we propose a control consistency metric. This metric works by estimating the control signal embedded in the generated sample and comparing it to the original input signal. Our results further corroborate the strong controllability achieved by OFT.
\end{itemize}

\section{Related Work}

\textbf{Text-to-image diffusion models}. The landscape of text-to-image synthesis has experienced remarkable advancements~\cite{saharia2022photorealistic,ramesh2022hierarchical,rombach2022high,gu2022vector,nichol2022glide}, driven primarily by progress in diffusion models~\cite{ho2020denoising,song2021score,song2021denoising,dhariwal2021diffusion} and sophisticated vision-language representation techniques~\cite{su2020vlbert,lu2019vilbert,radford2021learning,wang2021simvlm,li2022blip,li2023blip,alayrac2022flamingo,schuhmann2022laion}. Pixel-space diffusion models, such as GLIDE~\cite{nichol2022glide} and Imagen~\cite{saharia2022photorealistic}, differ in their treatment of the text encoder: GLIDE jointly learns the text encoder with the diffusion process using text-image pairs, whereas Imagen employs a pretrained, fixed text encoder. In contrast, latent-space models like Stable Diffusion~\cite{rombach2022high} and DALL-E2~\cite{ramesh2022hierarchical} employ a different representation strategy. Stable Diffusion uses VQ-GAN~\cite{esser2021taming} for its latent codebook, and DALL-E2 utilizes CLIP~\cite{radford2021learning} for building a joint embedding space across image and text modalities. Beyond diffusion approaches, generative adversarial networks~\cite{reed2016generative,zhang2017stackgan,xu2018attngan,li2019controllable} and autoregressive frameworks~\cite{ramesh2021zero,ding2021cogview,wu2022nuwa,yuscaling} have also been explored for text-to-image applications. Importantly, OFT is finetuning method-agnostic and can be seamlessly applied to any text-to-image diffusion model.

\textbf{Subject-driven generation}. To minimize alteration of the subject during generation, techniques such as those in \cite{avrahami2022blended,nichol2022glide} incorporate user-specified masks as additional conditions. Inversion-based approaches~\cite{choi2021ilvr,dhariwal2021diffusion,ramesh2022hierarchical,gal2022image} modify contextual aspects while retaining the core subject, and methods such as \cite{hertz2022prompt} achieve localized or global editing without relying on explicit input masks. However, many of these fail to robustly preserve subject identity features. Pivotal Tuning~\cite{roich2022pivotal} addresses this by finetuning generators around an inverted code and applying regularization to maintain identity, while \cite{nitzan2022mystyle} derives a personal generative face prior from multiple images of an individual. Although \cite{casanova2021instance} can synthesize various instances of an object, it may overlook instance-specific details. DreamBooth~\cite{ruiz2023dreambooth} customizes the diffusion model using a unique token, several subject images, and losses for both reconstruction and class-specific prior preservation. In contrast, OFT builds upon the DreamBooth setup but departs from simple, low-rate finetuning by utilizing orthogonal transformations for parameter updates.

\textbf{Controllable generation}. Image-to-image translation is often framed as a controllable generation problem. Traditionally, many approaches employ conditional generative adversarial networks~\cite{isola2017image,zhu2017toward,wang2018high,park2019semantic,choi2018stargan} to address this task. More recently, diffusion models have also become a popular choice for image-to-image translation~\cite{saharia2022palette,voynov2022sketch,wang2022pretraining}. Among newer developments, ControlNet~\cite{zhang2023adding} presents a technique that steers a pretrained diffusion model via finetuning, thereby incorporating additional control signals and attaining strong controllable generation capabilities. Similarly, T2I-Adapter~\cite{mou2023t2i} finetunes pretrained diffusion models to enhance control over generated images. Adopting the evaluation protocols established by~\cite{zhang2023adding,mou2023t2i}, we utilize OFT to finetune pretrained diffusion models and demonstrate superior controllability with reduced training data and fewer tunable parameters. Importantly, OFT incurs no extra computational cost during test-time inference.

\textbf{Model finetuning}. The practice of finetuning large pretrained models for downstream applications has garnered significant attention in recent years~\cite{devlin2018bert,brown2020language,he2022masked}. Within this paradigm, adaptation strategies such as those described in~\cite{hulora2022,houlsby2019parameter,pfeiffer2020adapterfusion} have been thoroughly investigated in the context of natural language processing. OFT's closest counterpart is LoRA~\cite{hulora2022}, which employs a low-rank decomposition to update weights additively during finetuning. In contrast, OFT implements a shared orthogonal transformation across layers, updating neuron weights multiplicatively. This approach mathematically guarantees preservation of pairwise angular relationships among neurons within the same layer, thereby promoting greater stability.

\section{Orthogonal Finetuning}

\subsection{Why Does Orthogonal Transformation Make Sense?}

To motivate the application of orthogonal transformation for finetuning text-to-image diffusion models, we break down the reasoning into two main points: (1) the rationale for finetuning the direction (or angle) of neuron weights, and (2) the justification for selecting orthogonal transformations as the method to adjust these directions.

Addressing the first question, we take cues from prior empirical findings \cite{liu2018decoupled,chen2020angular} suggesting that differences in angular features serve as an effective metric for the semantic gap. SphereNet~\cite{liu2017deep} demonstrated that constraining neurons to a unit hypersphere during training maintains model capacity while often enhancing generalization, which suggests that neuron direction encodes critical data characteristics. To elucidate the role of neuron orientation, we present a toy experiment (see Figure~\ref{fig:toy_ae}) utilizing a basic convolutional autoencoder trained on flower images. During training, feature maps are generated via standard inner products between the convolution kernel $\bm{w}$ and sliding-window input $\bm{x}$ (with $z$ as the resulting element). In the evaluation phase, we assess three feature map constructions: (a) retaining the training-phase inner product, (b) isolating magnitude, and (c) isolating angular components. Results in Figure~\ref{fig:toy_ae} reveal that neuron angular data nearly perfectly reconstructs input images, whereas magnitude alone conveys negligible information. Notably, the angular (cosine) transformation (c) is excluded during training, which only involves the inner product. These findings underscore that neuron direction predominantly encodes the input’s semantic content, implying that semantic manipulation may be most effectively achieved through tuning neuron directions.

For the second question, the most straightforward method to adjust neuron directions is to apply uniform rotation or reflection to all neurons within a layer, naturally leading to orthogonal transformations. While transformations allowing distinct angular adjustments per neuron may increase adaptability, we observe that orthogonal transformations offer an optimal balance of flexibility and constraint. Supporting this, \cite{liu2021orthogonal} demonstrates the expressive power of orthogonal transformations in neural network learning. To substantiate our position, we perform a controlled generation experiment using the ControlNet~\cite{zhang2023adding} framework, as illustrated in Figure~\ref{fig:ortho}. Results show that our standard OFT exhibits stable performance and attains precise control post-training (by epoch 20). In contrast, omitting the orthogonality constraint from OFT leads to failure in producing convincing images or achieving control objectives. This experiment confirms the critical role of the orthogonality constraint in the effectiveness of OFT.

\subsection{General Framework}

Traditionally, finetuning involves employing gradient descent with a reduced learning rate to adjust a model's parameters, either wholly or for specific layers. This modest learning rate acts as a soft regularizer, discouraging substantial deviation from the pretrained weights; thus, standard finetuning inherently attempts to limit the change, effectively minimizing $\thickmuskip=2mu \medmuskip=2mu \| \bm{M} - \bm{M}^0\|$—with $\bm{M}$ representing the updated model weights and $\bm{M}^0$ the original pretrained ones. While this constraint is reasonable, it may still allow excessive flexibility, particularly with large-scale models. To tackle this, LoRA imposes an explicit low-rank structure on the weight adaptation, specifying $\thickmuskip=2mu \medmuskip=2mu \text{rank}(\bm{M} - \bm{M}^0)=r'$, where $r'$ is a predetermined small integer. Unlike LoRA, OFT (Orthogonal Finetuning) places a restriction based on pairwise neuron relationships by enforcing $\thickmuskip=2mu \medmuskip=2mu \| \text{HE}(\bm{M})-\text{HE}(\bm{M}^0) \|=0$, where the operator $\text{HE}(\cdot)$ computes the hyperspherical energy of the given weight matrix.

To illustrate, consider a fully connected layer $\thickmuskip=2mu \medmuskip=2mu \bm{W}=\{\bm{w}_1,\cdots,\bm{w}_n\}\in\mathbb{R}^{d\times n}$, where each column $\bm{w}_i\in\mathbb{R}^{d}$ denotes the $i$-th neuron's parameters, and $\bm{W}^0$ is the associated pretrained weight matrix. For an input vector $\thickmuskip=2mu \medmuskip=2mu \bm{x}\in\mathbb{R}^d$, the layer's output is $\thickmuskip=2mu \medmuskip=2mu \bm{z}=\bm{W}^\top\bm{x}$. Under the OFT paradigm, the goal is to keep the difference in hyperspherical energy between finetuned and original weights minimal:

\begin{equation}\label{eq:oft_principle}
\footnotesize
\min_{\bm{W}} \left\| \text{HE}(\bm{W})-\text{HE}(\bm{W}^0)\right\|~~\Leftrightarrow~~\min_{\bm{W}} \bigg{\|} \sum_{i\neq j} \|\hat{\bm{w}}_i-\hat{\bm{w}}_j\|^{-1}-\sum_{i\neq j} \|\hat{\bm{w}}^0_i-\hat{\bm{w}}^0_j\|^{-1}\bigg{\|}
\end{equation}

Here, $\thickmuskip=2mu \medmuskip=2mu \hat{\bm{w}}_i:=\bm{w}_i/\|\bm{w}_i\|$ refers to the unit-normalized $i$-th neuron, and the hyperspherical energy for the layer $\bm{W}$ is expressed as $\thickmuskip=2mu \medmuskip=2mu \text{HE}(\bm{W}):= \sum_{i\neq j} \|\hat{\bm{w}}_i-\hat{\bm{w}}_j\|^{-1}$. It is straightforward to see that this objective reaches a minimum value of zero in Eq.~\eqref{eq:oft_principle} when the weight matrices $\bm{W}$ and $\bm{W}^0$ only differ by an orthogonal transformation—i.e., $\thickmuskip=2mu \medmuskip=2mu \bm{W}=\bm{R}\bm{W}^0$ for some orthogonal $\bm{R}\in\mathbb{R}^{d\times d}$. Here, $\det(\bm{R})=1$ indicates a rotation and $\det(\bm{R})=-1$ indicates a reflection. The core idea of OFT, therefore, is to finetune the model exclusively via layer-wise orthogonal transformations, optimizing the shared orthogonal matrix per layer. More precisely, OFT seeks to determine the optimal orthogonal matrix $\thickmuskip=

\begin{equation}\label{eq:oft_general}
    \footnotesize
    \bm{z} = \bm{W}^\top \bm{x} = (\bm{R} \cdot \bm{W}^0)^\top \bm{x},~~~\text{s.t.}~\bm{R}^\top \bm{R} = \bm{R} \bm{R}^\top = \bm{I}
\end{equation}

In this formulation, $\bm{W}$ refers to the weight matrix updated via OFT, and $\bm{I}$ stands for the identity matrix. The structure of OFT is depicted in Figure~\ref{fig:block}. Analogous to the practice of zero initialization in LoRA, it is essential that OFT begins its fine-tuning precisely from the pretrained parameter $\bm{W}^0$. To facilitate this, we set the initial value of the orthogonal matrix $\bm{R}$ as the identity matrix, thereby ensuring that the optimization proceeds from the pretrained weights. To maintain the orthogonality constraint on $\bm{R}$, differential orthogonalization methods as outlined in \cite{liu2021orthogonal,lezcano2019cheap} can be employed. Later in the discussion, we elaborate on strategies to preserve orthogonality in a computationally economical manner.

\subsection{Efficient Orthogonal Parameterization}\label{sect:eff_ortho}

Conventional approaches for enforcing orthogonality, such as the Gram-Schmidt process, while differentiable, often impose significant computational overhead~\cite{liu2021orthogonal}. To improve efficiency, we utilize Cayley parameterization as a means to generate orthogonal matrices. More specifically, the orthogonal matrix $\bm{R}$ is expressed as $\thickmuskip=2mu \medmuskip=2mu \bm{R} = (\bm{I} + \bm{Q})(\bm{I} - \bm{Q})^{-1}$, where the matrix $\bm{Q}$ satisfies skew-symmetry, i.e., $\thickmuskip=2mu \medmuskip=2mu \bm{Q} = -\bm{Q}^\top$. This approach offers computational benefits, but with a constraint: Cayley parameterization restricts us to constructing only special orthogonal matrices (i.e., determinant $1$). Empirically, however, this restriction is negligible and does not degrade practical performance. Even so, the Cayley-transformed orthogonal matrix $\bm{R}$ can become storage-intensive for large $d$. To alleviate this issue, we propose constructing $\bm{R}$ as a block-diagonal matrix with $r$ constituent blocks, resulting in the following structure:

\begin{equation}
    \footnotesize
    \bm{R} = \text{diag}(\bm{R}_1, \bm{R}_2, \cdots, \bm{R}_r) = \begin{bmatrix}
    \bm{R}_1 \in \text{O}(\frac{d}{r}) &  &\\
    &\ddots & \\
    & & \bm{R}_r \in \text{O}(\frac{d}{r})
    \end{bmatrix} \in \text{O}(d)
\end{equation}

Here, $\text{O}(d)$ represents the group of orthogonal matrices in $d$ dimensions, with $\thickmuskip=2mu \medmuskip=2mu \bm{R}\in\mathbb{R}^{d\times d}$ and $\thickmuskip=2mu \medmuskip=2mu \bm{R}_i\in\mathbb{R}^{d/r\times d/r}$ for each $i$ being orthogonal. In the special case where $\thickmuskip=2mu \medmuskip=2mu r=1$, the block-diagonal orthogonal matrix reduces to an ordinary unconstrained orthogonal matrix. For a full orthogonal matrix of size $\thickmuskip=2mu \medmuskip=2mu d\times d$, the parameter count is $\thickmuskip=2mu \medmuskip=2mu d(d-1)/2$, thus incurring a complexity of $\mathcal{O}(d^2)$. In contrast, if a block-diagonal structure with $r$ blocks is employed, the parameter count decreases to $\thickmuskip=2mu \medmuskip=2mu d(d/r-1)/2$, yielding a complexity of $\thickmuskip=2mu \medmuskip=2mu \mathcal{O}(d^2/r)$. To reduce parameters even further, one may enforce block sharing, specifically, $\thickmuskip=2mu \medmuskip=2mu \bm{R}_i=\bm{R}_j$ for all $i\neq j$, leading to parameter complexity of $\mathcal{O}(d^2/r^2)$. Despite these parameter-saving techniques, the constructed matrix $\bm{R}$ retains orthogonality, ensuring the preservation of the hyperspherical property.

Let us now compare OFT and LoRA regarding their parameter requirements. For LoRA, with rank $r'$, the number of parameters optimized during training is given by $\thickmuskip=2mu \medmuskip=2mu r'(d+n)$. If we set both $r$ and $r'$ as fractions of the neuron dimension $d$ (for instance, $\thickmuskip=2mu \medmuskip=2mu r=r'=\alpha d$ with $\thickmuskip=2mu \medmuskip=2mu 0<\alpha\leq 1$), the parameter complexity for LoRA rises to $\mathcal{O}(d^2+dn)$, whereas for OFT it becomes $\mathcal{O}(d)$. To highlight this difference with a practical scenario, consider a $128\times128$ weight matrix: if $\thickmuskip=2mu \medmuskip=2mu r'=8$, LoRA would involve $2,048$ tunable parameters, while OFT with $\thickmuskip=2mu \medmuskip=2mu r=8$ (no block sharing) would require $960$.

\subsection{Constrained Orthogonal Finetuning}

One can tighten the original OFT by confining the finetuned model within a limited range of the pretrained parameters. This is the essence of COFT, where the model output is computed as:

\begin{equation}\label{eq:coft_general}
    \footnotesize
    \bm{z}=\bm{W}^\top\bm{x}=(\bm{R}\cdot\bm{W}^0)^\top\bm{x},~~~\text{s.t.}~\bm{R}^\top\bm{R}=\bm{R}\bm{R}^\top=\bm{I},~\norm{\bm{R}-\bm{I}}\leq \epsilon
\end{equation}

In this formulation, $\bm{R}$ is constrained to be orthogonal, with an additional restriction such that its deviation from the identity matrix is bounded by $\epsilon$. The orthogonality condition may be enforced using the Cayley parameterization from Section~\ref{sect:eff_ortho}. Nonetheless, imposing the $\epsilon$-deviation constraint directly within the Cayley-parameterized structure is challenging. To develop a better understanding of the Cayley transform, we can use a Neumann series expansion to approximate $\thickmuskip=2mu \medmuskip=2mu \bm{R}=(\bm{I}+\bm{Q})(I-\bm{Q})^{-1}$ by $\thickmuskip=2mu \medmuskip=2mu

Since the series converges in the operator norm, the constraint $\thickmuskip=2mu \medmuskip=2mu\|\bm{R}-\bm{I}\|\leq\epsilon$ may be equivalently reformulated by mapping it into the domain of the Cayley transform. This produces the alternative constraint $\thickmuskip=2mu \medmuskip=2mu \|\bm{Q}-\bm{0}\|\leq\epsilon'$, where $\bm{0}$ represents the zero matrix and $\epsilon'$ is a hyperparameter specifying the permitted deviation (distinct from $\epsilon$). The resulting condition on $\bm{Q}$ is conveniently handled using projected gradient descent. To ensure that the orthogonal matrix $\bm{R}$ is initialized as the identity, $\bm{Q}$ is set to the zero matrix during initialization. The COFT method thus implements two direct constraints: one minimizes the difference in hyperspherical energy, while the other explicitly restricts the deviation from the initial, pretrained model. While the second restriction is commonly applied in an implicit manner by prevailing finetuning approaches, COFT introduces it as an explicit design choice. Although OFT demonstrates strong performance, we have found that the explicit deviation constraint in COFT leads to further improvements in finetuning stability compared to OFT. The impact of the parameter $\epsilon$ on COFT is illustrated in Figure~\ref{fig:eps_coft}. The flexibility of finetuning under COFT is modulated by $\epsilon$: larger $\epsilon$ values make COFT’s results increasingly similar to those of OFT, while smaller $\epsilon$ yields behavior closer to the original pretrained text-to-image diffusion model.

\subsection{Re-scaled Orthogonal Finetuning}

We introduce a straightforward modification to standard OFT, where, in addition to the orthogonal transformation, a separate scaling factor is learned for each neuron. This enhancement is motivated by the observation that scaling neuron activations does not affect their hyperspherical energy, since normalization negates the magnitude. Specifically, the forward propagation is performed as $\thickmuskip=2mu \medmuskip=2mu\bm{z}=(\bm{R}\bm{W}^0\bm{D})^\top\bm{x}$\footnote{\textbf{Errata}: The NeurIPS camera-ready version incorrectly lists the forward computation of re-scaled OFT as $\thickmuskip=2mu \medmuskip=2mu\bm{z}=(\bm{D}\bm{R}\bm{W}^0)^\top\bm{x}$. However, the actual implementation is correct, thus all results in Appendix~\ref{app:rescaled} are valid.}, where $\thickmuskip=2mu \medmuskip=2mu \bm{D}=\text{diag}(s_1,\cdots,s_n)\in\mathbb{R}^{n\times n}$ is a learnable diagonal matrix whose entries $s_1,\ldots,s_n$ are strictly positive. Unlike the original OFT, which restricts learning to $\bm{R}$, this variation permits both $\bm{D}$ and $\bm{R}$ to be updated during training. Introducing re-scaling modestly expands the parameter count but confers significantly greater flexibility. For our main experimental analysis, we adhere to the standard OFT framework to isolate the influence of orthogonal transformations; however, in practice, re-scaled OFT generally offers superior results (refer to Appendix~\ref{app:rescaled}).

\section{Intriguing Insights and Discussions}

\textbf{OFT is broadly compatible with various architectures}. In theory, OFT can be implemented in any neural architecture. For instance, while LoRA is predominantly used on attention matrices within Transformers~\cite{hulora2022}, for parity in our evaluations, we restrict OFT to finetuning attention weights in these models. Additionally, OFT’s design also renders it highly effective for convolutional layers. Here, the block-diagonal form of $\bm{R}$ lends itself to interesting interpretations, in contrast to LoRA’s parameterization. Should the number of blocks match the count of input channels, each block acts independently on a single neuron channel—mirroring the operation of depth-wise convolution kernels~\cite{chollet2017xception}. Conversely, if $\bm{R}$ is structured to share all blocks, it realizes an orthogonal transformation shared across channels. For preliminary experiments applying OFT to convolutional finetuning, please refer to Appendix~\ref{app:convolution}.

\textbf{Connection to LoRA}. LoRA introduces a low-rank matrix to ensure that modifications to the pretrained weight matrix do not drastically alter its information content. In comparison, OFT employs an orthogonal (and thus full-rank) transformation to regulate changes applied to the pretrained weights, thereby preserving the integrity of the pretrained information. The OFT forward computation can be reformulated as $\thickmuskip=2mu \medmuskip=2mu \bm{z}=(\bm{R}\bm{W}^0)^\top\bm{x}=(\bm{W}^0+(\bm{R}-\bm{I})\bm{W}^0)^\top\bm{x}$, where $\thickmuskip=2mu \medmuskip=2mu (\bm{R}-\bm{I})\bm{W}^0$ serves a similar role to LoRA’s low-rank update. Given that $\bm{W}^0$ is generally of full rank, OFT achieves a low-rank modification as well when $\thickmuskip=2mu \medmuskip=2mu \bm{R}-\bm{I}$ itself is low-rank. In analogy with LoRA’s rank hyperparameter $r'$, OFT uses a diagonal block size parameter $r$ to restrict the number of parameters that require training. Notably, while both LoRA and OFT seek parameter efficiency, their approaches are distinct: LoRA utilizes low-rank structure to economize parameters, whereas OFT leverages block-diagonal sparsity by enforcing orthogonality.

\textbf{Why OFT converges faster}. One explanation, as illustrated in Figure~\ref{fig:toy_ae}, is that modifying the orientation of neurons—OFT’s explicit goal—is the most direct way to adapt semantic representations. Moreover, OFT’s approach can be understood as finetuning the neurons while constraining them to move on a smooth hypersphere, which contributes to a more favorable optimization landscape. This smoother landscape and improved convergence are further corroborated in \cite{liu2021orthogonal}.

\textbf{Why not minimize hyperspherical energy}. Unlike \cite{liu2021orthogonal}, this work does not pursue minimization of hyperspherical energy. For classification tasks, redundancy-free neurons are optimal, yet minimizing hyperspherical energy enforces a uniform spacing of all neurons across the hypersphere. Such a criterion is counterproductive during finetuning, as it risks erasing information accrued during pretraining.

\textbf{Balancing Flexibility and Regularity in Finetuning}. Our analysis reveals a fundamental trade-off between flexibility and regularity. Although standard finetuning provides maximal flexibility, it is prone to instability and can frequently result in model collapse. In contrast, OFT, despite its straightforward design, achieves a robust compromise between these two aspects by maintaining the pairwise angular relationships among neurons. The block-diagonal parameterization further enforces regularization by imposing stricter constraints on the orthogonal transformation.

\textbf{Inference Efficiency Unaffected}. In distinction to ControlNet, our OFT method incurs no inference-time overhead in the finetuned network. During inference, the learned orthogonal transformation $\bm{R}$ can be directly applied to the pretrained weight matrix $\bm{W}^0$, yielding the final weight matrix $\thickmuskip=2mu \medmuskip=2mu \bm{W}=\bm{R}\bm{W}^0$. As a result, inference proceeds at the same speed as with the original pretrained weights.

\section{Experiments and Results}

\textbf{Experimental Protocols}. For our evaluation, we utilize Stable Diffusion v1.5~\cite{rombach2022high} as the pretrained text-to-image foundation. To maintain fairness, we sample generated outputs from all compared approaches randomly. In scenarios involving subject-driven image synthesis, we broadly adhere to DreamBooth~\cite{ruiz2023dreambooth}, while for controllable generation tasks, our setup aligns with methods such as ControlNet~\cite{zhang2023adding} and T2I-Adapter~\cite{mou2023t2i}. To enable an equitable comparison with LoRA, we restrict OFT and COFT application to identical layers as those modified by LoRA. Supplementary results and methodological particulars are detailed in Appendix~\ref{app:settings}.

\subsection{Subject-driven Generation}

\textbf{Implementation Details}. DreamBooth~\cite{ruiz2023dreambooth} and LoRA~\cite{hulora2022} serve as our primary baselines. Each method is trained using the same objective function as prescribed by DreamBooth. Where applicable, the original hyperparameters and protocol recommendations from the respective papers are employed for DreamBooth and LoRA. Further experimental data are provided in Appendix~\ref{app:settings},\ref{app:coft_oft},\ref{app:more_qualitative},\ref{app:failure}.

\textbf{Finetuning stability and convergence}. To begin, we assess both the stability and rate of convergence during finetuning for DreamBooth, LoRA, OFT, and COFT. The corresponding findings are displayed in Figure~\ref{fig:teaser} and Figure~\ref{fig:db_convergence}. From these visualizations, it is clear that COFT and OFT facilitate stable finetuning of the diffusion model. After 400 training steps, effective control is evident in both DreamBooth and OFT variants, whereas LoRA fails to accurately retain the subject’s identity. Progressing to 2000 iterations, DreamBooth starts producing degraded (collapsed) outputs, and LoRA ceases to render a yellow shirt, producing yellow fur instead. In contrast, OFT and COFT maintain reliable and stable control over the generated content throughout. This evidence highlights the rapid and robust convergence properties of the OFT framework in the context of subject-driven image synthesis. It is important to emphasize the robustness with respect to the number of finetuning iterations: this feature lessens the burden of manually adjusting the iteration count across different subjects. For OFT and COFT, one can use a relatively high number of iterations without meticulous calibration. In addition, with a suitable choice of $\epsilon$ for COFT, both the learning rate and the number of training steps are straightforward to configure.

\textbf{Quantitative comparison}. Building on the evaluation approach in \cite{ruiz2023dreambooth}, we present quantitative comparisons regarding subject identity preservation (using DINO~\cite{caron2021emerging}, CLIP-I~\cite{radford2021learning}), adherence to the text prompt (CLIP-T~\cite{radford2021learning}), and the diversity of generated samples (LPIPS~\cite{zhang2018unreasonable}). In detail, CLIP-I quantifies the mean pairwise cosine similarity between CLIP feature embeddings for real versus synthesized images, while DINO employs ViT S/16 DINO embeddings in a similar setup. For text alignment, CLIP-T evaluates the mean cosine similarity between the embeddings of the input prompt and generated images. The diversity measure consists of calculating the mean LPIPS cosine similarity among images of the same subject generated with the same text input. As summarized in Table~\ref{table:dreambooth_final}, both COFT and OFT substantially surpass DreamBooth and LoRA in DINO and CLIP-I scores, and provide at least comparable—often superior—performance in terms of prompt adherence and diversity. To account for stochastic variability, each method is finetuned thirty times on the same pretrained model using distinct random seeds. Collectively, the results indicate that OFT not only ensures improved convergence stability, but also consistently delivers enhanced final results.

\textbf{Qualitative comparison}. To provide clearer insight into the advantages of OFT, we present a selection of randomly chosen examples showcasing subject-driven generation in Figure~\ref{fig:dreambooth_final}. For consistency, all samples are produced using the same finetuned model with each respective technique, and no separate prompt optimization is carried out for any result. For every method, we report outcomes from the model yielding the highest validation CLIP scores. Examination of Figure~\ref{fig:dreambooth_final} reveals that both OFT and COFT maintain strong semantic fidelity to the subject, whereas LoRA frequently struggles with subject identity retention (for instance, LoRA completely fails in the bowl example). Additionally, OFT and COFT display substantially finer control over generation according to textual prompts, whereas DreamBooth, though effective at subject identity preservation, often fails to align image outputs with the prompt (as seen in the first row for the bowl example). This comparison highlights that the OFT framework attains superior balance between controllability and subject fidelity. Furthermore, OFT demonstrates robust performance across varying iteration numbers, in contrast to both DreamBooth and LoRA, which are more iteration-sensitive. Expanded qualitative results are provided in Appendix~\ref{app:more_qualitative}, while further evidence of OFT's effectiveness is supported by a human study discussed in Appendix~\ref{app:human}.

\subsection{Controllable Generation}

\textbf{Settings}. As comparative baselines, we utilize ControlNet~\cite{zhang2023adding}, T2I-Adapter~\cite{mou2023t2i}, and LoRA~\cite{hulora2022}. Our primary evaluations consider three complex controllable generation scenarios: converting Canny edge maps to images~(C2I) using the COCO dataset~\cite{lin2014microsoft}, transforming segmentation maps to images~(S2I) on ADE20K~\cite{zhou2017scene}, and generating faces from landmarks~(L2F) using the CelebA-HQ dataset~\cite{karras2018progressive,xia2021tedigan}. Each method is used to finetune Stable Diffusion~(SD) v1.5 for 20 epochs on these tasks. Additional qualitative findings can be found in Appendix~\ref{app:more_qualitative},\ref{app:more_control_task},\ref{app:failure}.

\textbf{Convergence}. We assess how rapidly ControlNet, T2I-Adapter, LoRA, and COFT converge on the L2F task using both quantitative and qualitative analyses. As our quantitative metric, we calculate the mean $\ell_2$ distance between the set of control face landmarks and the corresponding predicted face landmarks. Figure~\ref{fig:face_convergence} depicts the error in landmark localization as a function of the number of fine-tuning epochs for each model. It is clear that COFT and OFT reach low error values considerably sooner than their counterparts. For instance, LoRA requires approximately 20 epochs to attain the same accuracy level that OFT achieves after just 8 epochs. Importantly, OFT and COFT employ a number of trainable parameters on par with LoRA (and significantly fewer than those in ControlNet), yet they deliver markedly swifter convergence compared to prevailing techniques. Additional evidence of OFT’s efficient learning is found in Figure~\ref{fig:teaser}, where the example on the right demonstrates OFT’s capacity to converge using as little as 5\% of the ADE20K dataset, in stark contrast to ControlNet and LoRA, which need substantially more data to reach comparable performance. Regarding qualitative comparisons, we concentrate on OFT, COFT, and ControlNet, as ControlNet exhibits a similar degree of landmark error. Figure~\ref{fig:face_convergence_qual} reveals that OFT and COFT display consistently stable and progressive alignment of the generated face pose with the control landmarks. In contrast, ControlNet shows a sudden improvement in controllability after epoch 8, consistent with the abrupt convergence observed by \cite{zhang2023adding}. The stability and predictability in the convergence dynamics of our OFT framework facilitate practical deployment, as the smooth progression of training makes the selection of optimal hyperparameters for OFT far more straightforward.

\textbf{Quantitative comparison}. To assess the effectiveness of controllable generation, we propose a metric for control consistency. This metric operates by extracting the control signal from the generated output and directly measuring its alignment with the input control signal. For the C2I task, we report IoU and F1 scores, while for S2I, we present mean IoU alongside mean and overall accuracy. The L2F task evaluation involves calculating the mean $\ell_2$ distance between the target and predicted control landmarks. Additional information on these evaluation metrics is provided in Appendix~\ref{app:settings}. All baseline methods are evaluated under their optimal hyperparameter configurations. As displayed in Table~\ref{table:control_gen}, OFT and COFT outperform competing approaches, achieving superior control fidelity and accuracy. In contrast, adapter-based models (\emph{e.g.}, T2I-Adapter, ControlNet) demonstrate slower convergence and suboptimal results. LoRA exhibits stronger performance than ControlNet for S2I but lags behind in C2I and L2F tasks. Overall, LoRA's upper performance bound remains limited despite extensive hyperparameter optimization. Notably, the OFT framework continues to show performance gains as the parameter count increases, suggesting its full potential is not yet reached. Importantly, our quantitative approach to evaluating controllable generation constitutes a key novelty, offering a precise means to measure control in finetuned models, subject only to the accuracy constraints of the underlying segmentation or detection models.

\textbf{Qualitative comparison}. We further conduct a qualitative evaluation of OFT and COFT against leading approaches such as ControlNet, T2I-Adapter, and LoRA. As depicted in the randomly selected examples in Figure~\ref{fig:control_final}, both OFT and COFT are capable of generating images with superior visual realism and faithful detail, while also demonstrating precise control over the generation process. In the S2I scenario, LoRA is unable to produce images consistent with the provided segmentation maps, whereas ControlNet, OFT, and COFT are effective at maintaining structural alignment. OFT and COFT, in particular, deliver visually compelling results with richer detail and more plausible geometry than ControlNet, despite employing considerably fewer parameters. Regarding the C2I task, OFT and COFT are successful in synthesizing images that capture the semantics of coarse Canny edge maps, while T2I-Adapter and LoRA are noticeably less effective. For the L2F task, our approach enables the most precise reproduction of target facial poses, even in difficult cases with extreme pose variation. Across all three categories of control tasks, both OFT and COFT consistently outperform current state-of-the-art baselines in qualitative image quality, which underlines the strong controllable generation capabilities offered by the OFT framework. For broader assessment, Appendix~\ref{app:control_exp} presents additional qualitative results across the three tasks, and Appendix~\ref{app:more_control_task} further illustrates the versatility of OFT on diverse tasks, such as dense pose to human body, sketch-to-image, and depth-to-image translation.

\section{Concluding Remarks and Open Problems}

Driven by the insight that angular relationships between neurons play a pivotal role in encoding visual semantics, we introduce a straightforward yet powerful fine-tuning technique—orthogonal finetuning (OFT)—for exerting control over text-to-image diffusion models. Our approach primarily focuses on two major text-to-image use cases: subject-driven generation and controllable generation. Relative to prior methodologies, OFT achieves enhanced controllability and more stable fine-tuning, all while utilizing a smaller number of tunable parameters. A key advantage of OFT is that it does not increase inference time, thereby facilitating efficient and practical deployment.

However, OFT presents several intriguing open questions. First, its enforcement of orthogonality is accomplished through Cayley parametrization, which requires computing a matrix inverse. This necessity can restrict OFT's scalability. While our use of block diagonal parametrization mitigates this to some extent, discovering methods for differentiable and accelerated matrix inversion remains an unresolved issue. Second, OFT exhibits promising compositionality: the orthogonal matrices generated from multiple OFT fine-tuning processes can be composed via multiplication, yielding a new orthogonal matrix. Whether the resulting set of orthogonal matrices effectively retains the knowledge from each downstream task is a compelling topic for further investigation. Lastly, the parameter efficiency offered by OFT hinges on the block diagonal design, an aspect that introduces unavoidable biases and diminishes flexibility. Identifying strategies for enhancing parameter efficiency without such biases stands as an important and open challenge.